\documentclass{article}
\usepackage[T1]{fontenc}
\usepackage{lmodern}
\usepackage{graphicx}
\usepackage{amsmath,amssymb}
\usepackage[a4paper,margin=2cm]{geometry}
\usepackage[absolute,overlay]{textpos}
\setlength{\TPHorizModule}{1mm}
\setlength{\TPVertModule}{1mm}
\usepackage[indonesian]{babel}
\usepackage{hyperref}
\setlength{\emergencystretch}{2em}
% unit: Halaman 15

\begin{document}

% === Halaman 15 (python/native) ===

Contoh 3: Misalkan barisan superincreasing adalah \{2, 3, 6, 13, 27, 52\}, dan 

m = 105, dan n = 31.

 Barisan non-superincreasing (atau normal) knapsack dihitung sbb:


2  31 mod 105 = 62


3  31 mod 105 = 93


6  31 mod 105 = 81


13  31 mod 105 = 88


27  31 mod 105 = 102


52  31 mod 105 = 37

 Jadi, kunci publik adalah \{62, 93, 81, 88, 102, 37\}, sedangkan kunci privat 

adalah \{2, 3, 6, 13, 27, 52\}.

15

\end{document}
